\subsection*{A second field with the real-zero obstruction removed}
Take the Hilbert class field of $\mathbb Q(\sqrt{-71})$. Reduction of positive definite binary quadratic forms gives seven classes. Representatives up to inversion are
\[
 Q_0=(1,1,18),\quad Q_1=(2,1,9),\quad Q_2=(4,3,5),\quad Q_3=(3,1,6).
\]
Indeed a reduced form has $a\le\sqrt{71/3}<5$; checking $|b|\le a$ gives the displayed classes and their inverses. The group of order seven is cyclic. A prime above $2$ generates it; the square of that prime ideal has norm $4$ and nontrivial class, and the only nonprincipal inverse pair representing $4$ is $Q_2$. This identifies the displayed indices with the respective powers of a generator up to inversion. Complex conjugation gives the dihedral group of order $14$ over $\mathbb Q$.

Put $\Theta_k(iy)=\sum_{x,v\in\mathbb Z}\exp(-2\pi yQ_k(x,v))$, and $q=e^{-2\pi y}$. Grouping integral ideals by their ideal class gives the weight-one Hecke forms
\[
 f_j(iy)=\tfrac12\Theta_0(iy)+
   \sum_{k=1}^3\cos(2\pi jk/7)\Theta_k(iy),\qquad 1\le j\le3.
\]
Here the two units account for the factor $1/2$ before inverse classes are grouped. The constants cancel. For $y\ge1/\sqrt{71}$ one has $q<12/25$; for example $2\pi/\sqrt{71}>37/50$ and the exponential Taylor polynomial of degree four at $37/50$ exceeds $25/12$. On $0<q\le12/25$, respectively, we obtain the uniform lower bounds
\[
             f_1(iy)/q\ge0.999,\qquad f_2(iy)/q\ge0.877,\qquad f_3(iy)/q\ge0.324.
\]
All three bounds are certified by \texttt{check\_level71.py} using rational arithmetic. In detail, let $r_k(n)$ count the representations of $n$ by $Q_k$ and $t=2\cos(2\pi j/7)$. Twice the $n$th coefficient is
$r_0(n)+t r_1(n)+(t^2-2)r_2(n)+(1-t-t^2)r_3(n)$.
The three roots of $t^3+t^2-2t-1$ lie respectively in
$[1.2469,1.2470]$, $[-0.4451,-0.4450]$ and $[-1.8020,-1.8019]$ by sign changes. Enumerate the coefficients up to $N=24$ by
$4aQ(x,v)=(2ax+bv)^2+71v^2$. The coefficients have absolute value at most the number of ideals of norm $n$, at most the divisor count and hence at most $n$. With $r=12/25$, the remaining normalized series has absolute value at most $r^N((N+1)/(1-r)+r/(1-r)^2)$. Interval Horner evaluation on the 24 consecutive intervals of width $1/50$ gives the displayed bounds after subtracting this tail.

Poisson summation gives $\Theta_k(i/(71y))=\sqrt{71}\,y\Theta_k(iy)$: the adjugate of each binary Gram matrix is integrally equivalent to the original matrix by a quarter turn. Thus every $f_j(iy)$ is positive on the whole imaginary axis. Its Mellin integral is positive at every real argument and converges there by rapid decay at both ends. In particular, the three completed $L$-functions are positive throughout the real axis. The one-sided alternative of Section~\ref{sec:rotation} is therefore available: for any one of the three nonprincipal inverse pairs the bound $S_a(x)\le x/14+O_\epsilon(x^{1/2+\epsilon})$ is equivalent, under Dirichlet GRH, to the critical-line conclusion for all three nonlinear factors. The real-axis positivity and its rational certificate are unconditional. The required bound for the prime sum is open.
